\section{Attention and Awareness}

The relationship between attention and awareness represents a crucial domain for understanding how consciousness emerges from patterns of energetic coherence \citep{vanBoxtel2010}. While traditional approaches often treat attention as a spotlight illuminating contents for conscious awareness, ECC suggests a more nuanced relationship where attention shapes patterns of coherence that enable or constrain conscious experience.

From the perspective of energetic coherent computation, attention operates not merely as a selection mechanism but as an active process that modulates the stability and integration of coherent energy states across neural populations \citep{pinotsis2023cytoelectric}. This modulation occurs through multiple mechanisms, from altered metabolic dynamics to modified patterns of synchronization between brain regions. The framework helps explain both how attention enables conscious awareness and why certain stimuli can influence behavior without entering awareness.

The role of astrocytic networks proves particularly significant for understanding the attention-awareness relationship \citep{pellerin2004neuroenergetics}. Through their extensive gap junction networks, astrocytes help maintain broader patterns of energetic coherence that support sustained attention while enabling rapid shifts in conscious focus. This biological implementation explains how attention can simultaneously enhance processing of selected information while maintaining awareness of broader contextual states.

Consider how attention modifies neural light cones - the spatiotemporal boundaries within which conscious integration can occur \citep{hales2022electromagnetism}. Focused attention appears to enhance local coherence within specific neural populations while potentially constraining the broader field of conscious integration. This creates what neuroscientists have observed as enhanced processing of attended stimuli at the cost of reduced awareness of unattended information. Rather than representing a simple trade-off, however, this reflects sophisticated management of coherent energy states across multiple scales.

The distinction between overt and covert attention gains new significance through ECC's framework \citep{hunt2024electromagnetic}. Overt attention, involving physical orientation toward stimuli, creates direct changes in sensory input and processing. Covert attention, by contrast, operates through purely internal modulation of coherent states. Both forms, however, work by establishing specific patterns of energetic coherence that shape conscious experience. This explains both their distinct characteristics and their capacity to operate simultaneously in coordinated ways.

The framework particularly illuminates what neuroscientists term pre-attentive processing - the rapid, parallel evaluation of sensory input that occurs before focused attention. Rather than representing a separate processing stage, ECC suggests how pre-attentive mechanisms emerge from broader patterns of coherent energy distribution that enable rapid detection of potentially significant stimuli. This provides a physical basis for understanding how attention can be automatically captured by salient events while maintaining voluntary control over sustained focus.

The relationship between attention and perceptual awareness becomes particularly complex when examining cases of attentional capture and inattentional blindness \citep{block1995confusion}. Through ECC's framework, attentional capture occurs when specific patterns of energetic coherence achieve sufficient stability to redirect conscious awareness, even against voluntary control. This helps explain why certain stimuli - like sudden movements or personally relevant information - can automatically draw attention while others require deliberate focus to enter awareness.

Inattentional blindness, where observers fail to notice clearly visible but unexpected stimuli, reveals important constraints on how patterns of coherent energy can be maintained across neural populations. Rather than representing a failure of sensory processing, such blindness emerges from the brain's limited capacity to maintain multiple independent patterns of coherence simultaneously. The famous "gorilla experiment" by Simons and Chabris demonstrates how focused attention on one task can prevent the establishment of coherent states necessary for conscious awareness of other significant events.

The phenomenon of divided attention further illuminates how conscious awareness emerges from patterns of energetic coherence \citep{zheng2025unbearable}. While people can apparently divide attention between multiple tasks, careful experimental work reveals that this often involves rapid switching rather than truly parallel processing. Through ECC, this limitation reflects fundamental constraints on maintaining separate patterns of coherence rather than simple computational capacity limits. However, the framework also explains how certain combinations of tasks can be genuinely performed in parallel when they establish complementary rather than competing patterns of coherence.

Clinical conditions affecting attention provide additional insight into the relationship between coherent energy states and conscious awareness \citep{bosl2025dynamical}. Disorders like ADHD appear to involve disruption of mechanisms that maintain stable patterns of coherence across time, while conditions like spatial neglect suggest how damage to specific neural regions can prevent the establishment of coherent states necessary for conscious awareness in particular domains. The framework helps explain both why such conditions manifest as they do and why certain therapeutic approaches prove effective.

The role of neuromodulators in attention gains new significance through ECC's emphasis on energetic coherence \citep{pinotsis2023vivo}. Rather than simply enhancing or suppressing neural activity, neurotransmitters like acetylcholine and norepinephrine appear to modify how patterns of coherence can be established and maintained across neural populations. This helps explain both how pharmaceutical interventions affect attention and awareness, and why their effects often prove complex and context-dependent.

The distinction between endogenous and exogenous attention parallels important differences in how patterns of coherence are established and maintained. Endogenous attention involves top-down modulation of coherent states based on internal goals, while exogenous attention reflects bottom-up capture by stimuli that automatically establish strong patterns of coherence. Rather than representing competing processes, these mechanisms work together to optimize conscious awareness for both planned behavior and rapid response to unexpected events.

The temporal dynamics of attention and awareness reveal important properties of how conscious states emerge from patterns of energetic coherence \citep{amil2024theta}. The well-documented attentional blink phenomenon, where processing one target impairs awareness of a second target presented shortly afterward, reflects constraints on how quickly new patterns of coherence can be established while maintaining stability. Rather than representing mere resource limitations, this demonstrates how consciousness requires specific temporal windows for integrating information into stable coherent states.

ECC provides novel insight into what neuroscientists term the neural correlates of consciousness and their relationship to attention \citep{seth2024conscious}. Where traditional approaches often struggle to distinguish neural activity supporting attention from that enabling conscious awareness, the framework suggests how these processes interact through different aspects of energetic coherence. Attention appears to modulate local patterns of coherence that can then achieve sufficient stability to enter conscious awareness, explaining both their close relationship and their potential dissociation in certain conditions.

The framework proves particularly valuable for understanding meditation and other practices that deliberately modify the relationship between attention and awareness \citep{da2024mathematical}. Rather than treating these as purely psychological techniques, ECC suggests how such practices work by establishing specific patterns of energetic coherence that enable broader or more nuanced conscious awareness. This explains both the immediate effects of meditation on perception and attention, and how sustained practice can lead to lasting changes in conscious experience.

Social attention - our capacity to track and respond to others' attentional states - takes on new significance through ECC's analysis \citep{butlin2023consciousnessartificialintelligenceinsights}. The ability to maintain coherent representations of where others are attending appears to emerge from specialized neural systems that have evolved to support social coordination. This suggests how shared attention serves as a fundamental mechanism for establishing coherent states across individuals, enabling the sophisticated social coordination characteristic of human societies.

Looking toward future research directions, ECC suggests several promising avenues for investigating attention and awareness \citep{van2025processes}. These include examining how different frequencies of neural oscillation contribute to maintaining coherent states, studying how attention modulates energy consumption across brain regions, and investigating how various forms of meditation might enhance the brain's capacity for maintaining multiple stable patterns of coherence simultaneously. Such research could help develop more effective interventions for attention-related disorders while deepening our understanding of consciousness itself.

The complex relationship between attention and awareness thus serves as a crucial test case for ECC's broader framework \citep{dennett1992time}. By showing how these processes emerge from and remain grounded in patterns of energetic coherence, while maintaining distinct functional roles, the theory provides new ways to understand both normal cognitive function and its disruption in various conditions. This understanding suggests new approaches to enhancing attention and awareness through techniques that work with rather than against the brain's fundamental organizing principles.
